\documentclass[11pt,a4paper]{article}

\usepackage[utf8]{inputenc}
\usepackage[T1]{fontenc}
\usepackage[margin=1in]{geometry}
\usepackage{graphicx}
\usepackage{listings}
\usepackage{xcolor}
\usepackage{hyperref}
\usepackage{enumitem}
\usepackage{booktabs}

\definecolor{codebg}{rgb}{0.96,0.96,0.96}
\lstdefinestyle{java}{
    backgroundcolor=\color{codebg},
    basicstyle=\ttfamily\footnotesize,
    keywordstyle=\color{blue!70!black}\bfseries,
    commentstyle=\color{gray},
    stringstyle=\color{purple},
    numbers=left, numberstyle=\tiny\color{gray},
    breaklines=true, frame=single, captionpos=b,
    showstringspaces=false, language=Java
}
\lstset{style=java}

\title{\textbf{PetCare}\\Smart Veterinary Clinic \& Pet Health Management System\\\large A DSA Project Report}
\author{Sai Nikhil \\ \texttt{25f2005507@ds.study.iitm.ac.in}}
\date{\today}

\begin{document}
\maketitle
\hrule
\vspace{1em}

\section{Introduction}
PetCare is a Java-based command-line application that models a veterinary
clinic and pet healthcare management platform. It manages pet profiles and
medical histories, veterinary appointments and treatments, vaccination
schedules and reminders, and clinic analytics records. The platform uses
classical data structures and algorithms to organise pet health records
efficiently and to optimise the underlying clinic services.

\section{Project Architecture}
The project is structured as a standard Maven module organised by concern:
\begin{itemize}[noitemsep]
  \item \textbf{petcare.models} -- domain entities (\texttt{Pet},
        \texttt{Treatment}, \texttt{Appointment}, \texttt{Vaccination},
        \texttt{Clinic}).
  \item \textbf{petcare.utils} -- \texttt{InputUtils},
        \texttt{DataGenerator}, \texttt{DisplayUtils} for shared concerns.
  \item \textbf{petcare.menu} -- a top-level \texttt{MainMenu} and one
        sub-menu per module for a clean CLI workflow.
  \item \textbf{petcare.methods} -- the algorithm implementations,
        grouped into \texttt{trees}, \texttt{graphs}, \texttt{sorting} and
        \texttt{algorithms}.
\end{itemize}

\section{Module--Algorithm Mapping}
\begin{table}[h]
\centering
\begin{tabular}{@{}lll@{}}
\toprule
\textbf{Module} & \textbf{Course Outcome} & \textbf{Implementations} \\ \midrule
M1 & Trees \& Balanced Search    & BST, AVL, traversals \\
M2 & Multiway Trees \& Ranges    & B-Tree, B+ Tree, Segment, Fenwick \\
M3 & Graph Algorithms            & BFS, DFS, MST (Kruskal / Prim) \\
M4 & Shortest Paths              & Dijkstra, Bellman-Ford, Floyd-Warshall, Topological Sort \\
M5 & Sorting \& Ranking          & Merge, Quick, Heap, Radix / Counting \\
M6 & Greedy \& DP                & Activity Selection, Fractional Knapsack, 0/1 Knapsack, LIS \\ \bottomrule
\end{tabular}
\caption{Where each algorithm fits in the PetCare domain.}
\end{table}

\section{Module 1 --- Trees \& Balanced Search Structures}
\paragraph{Use case.} Pet records are indexed by \texttt{PetID}. A plain
Binary Search Tree (BST) gives \(O(\log n)\) operations on average, but
becomes \(O(n)\) on adversarial insertions. An AVL Tree restores the
\(O(\log n)\) guarantee by re-balancing after each modification, which is
why it is used as the production index over medical records. Tree
traversals produce sorted treatment histories.

\section{Module 2 --- Multiway Trees \& Range Queries}
\paragraph{Use case.} Clinic record IDs use a B-Tree of order \(t\) for
disk-friendly multiway indexing. Vaccination schedules are exposed through
a B+ Tree--style range index that answers ``which vaccines are due in day
range \([l,r]\)?'' efficiently. Segment Trees give range-sum queries over
daily visit counts; Fenwick Trees provide cumulative treatment statistics
with \(O(\log n)\) prefix sums.

\section{Module 3 --- Graph Algorithms for Healthcare Networks}
\paragraph{Use case.} Clinics are modelled as graph vertices, referral
links as weighted edges. BFS finds the \emph{nearest facilities} by
hop-count; DFS analyses workflow connectivity and counts components;
Kruskal and Prim produce the minimum spanning tree representing the
cheapest collaboration backbone among clinics.

\section{Module 4 --- Shortest Path Optimisation}
\paragraph{Use case.} Dijkstra computes the fastest route to every clinic
from a given source (single-source, non-negative weights). Bellman-Ford
copes with negative referral weights and detects negative-weight cycles.
Floyd-Warshall produces the full all-pairs distance matrix. Treatment
procedures (admit, diagnose, X-ray, surgery, recovery, discharge) form a
DAG, so a Kahn-style topological sort schedules them.

\section{Module 5 --- Advanced Sorting \& Ranking}
\paragraph{Use case.} Merge Sort sorts treatment records stably by cost.
Quick Sort ranks diagnoses by frequency in descending order. Heap Sort
extracts the top-\(K\) most common pet diseases. Radix and Counting Sort
sort pet identification numbers in linear time when keys are bounded.

\section{Module 6 --- Greedy \& Dynamic Programming}
\paragraph{Use case.} Greedy Activity Selection schedules the maximum
number of non-overlapping appointments. Fractional Knapsack allocates
divisible clinic resources within capacity. 0/1 Knapsack picks an optimal
combination of treatment plans within budget. Longest Increasing
Subsequence (patience-sort \(O(n \log n)\) variant) detects the longest
streak of improvement in a pet's weekly health score.

\section{How to Run}
\begin{lstlisting}[language=bash,frame=single]
# In STS / Eclipse:
File -> Import -> Maven -> Existing Maven Projects -> select PetCare

# From terminal:
javac -d out $(find src/main/java -name "*.java")
java -cp out petcare.PetCareApp
\end{lstlisting}

\section{Conclusion}
PetCare gathers six modules of foundational DSA into a single coherent
application. Each algorithm has a meaningful place in the veterinary
domain rather than being demonstrated in isolation, which makes the
project useful both as a reference implementation and as a teaching
exercise.

\end{document}
